\subtitle{\bf Aula 15 - Tipos em Linguagens de Programaçãoo}
\begin{frame}[t,plain]
\titlepage
\end{frame}

\begin{frame}{Objetivos}
\protect\hypertarget{objetivos-1}{}
\begin{itemize}
\item
  Apresentar sistemas de tipos como a especificação da análise
  semântica.
\item
  Discutir propriedades importantes de sistemas de tipos.
\end{itemize}
\end{frame}

\hypertarget{tipos-em-linguagens-de-programauxe7uxe3o}{%
\section{Tipos em Linguagens de
Programaçãoo}\label{tipos-em-linguagens-de-programauxe7uxe3o}}

\begin{frame}{O que é um Tipo?}
\protect\hypertarget{o-que-uxe9-um-tipo}{}
\begin{itemize}

\item
  Um \textbf{tipo} classifica valores e define:

  \begin{itemize}

  \item
    Quais operações são aplicáveis
  \item
    Como os valores são representados na memória
  \item
    Que conjunto de valores o tipo compreende
  \end{itemize}
\end{itemize}
\end{frame}

\begin{frame}[fragile]{Sistema de Tipo}
\protect\hypertarget{sistema-de-tipo}{}
\begin{itemize}

\item
  O \textbf{sistema de tipos} associa tipos a expressões e verifica a
  coerência dessas associações
\item
  Erros de tipo: somar um booleano a um inteiro, aplicar
  \texttt{pred true}
\item
  Expressões com erros de tipo são \textbf{sintaticamente válidas}, mas
  \textbf{semanticamente sem significado}
\end{itemize}
\end{frame}

\begin{frame}{Tipagem Estática}
\protect\hypertarget{tipagem-estuxe1tica}{}
\textbf{Tipagem estática} (\emph{static typing}): verificação em
\textbf{tempo de compila\~{c}ão}

\begin{itemize}

\item
  O compilador rejeita o programa se houver erro de tipo
\item
  Exemplos: Haskell, Java, C, Rust
\end{itemize}
\end{frame}

\begin{frame}{Tipagem Dinâmica}
\protect\hypertarget{tipagem-dinuxe2mica}{}
\begin{itemize}

\item
  \textbf{Tipagem dinâmica} (\emph{dynamic typing}): verificação em
  \textbf{tempo de execução}

  \begin{itemize}

  \item
    Erros detectados apenas quando a operação inválida é executada
  \item
    Exemplos: Python, JavaScript, Ruby, Lisp
  \end{itemize}
\end{itemize}
\end{frame}

\begin{frame}{Tipagem Forte}
\protect\hypertarget{tipagem-forte}{}
\begin{itemize}

\item
  \textbf{Tipagem forte} (\emph{strong typing}): não permite conversões
  implícitas entre tipos incompatíveis

  \begin{itemize}

  \item
    Operações entre tipos diferentes resultam em erro
  \item
    Exemplos: Haskell, Python
  \end{itemize}
\end{itemize}
\end{frame}

\begin{frame}[fragile]{Tipagem Fraca}
\protect\hypertarget{tipagem-fraca}{}
\begin{itemize}

\item
  \textbf{Tipagem fraca} (\emph{weak typing}): realiza conversões
  implícitas entre tipos

  \begin{itemize}

  \item
    Comportamentos frequentemente surpreendentes
  \item
    C: converte implicitamente entre inteiros e ponteiros
  \item
    JavaScript: \texttt{1 + "1"} resulta em
    \texttt{"11"}
  \end{itemize}
\end{itemize}
\end{frame}

\begin{frame}{Fraca vs Forte}
\protect\hypertarget{fraca-vs-forte}{}
\begin{itemize}

\item
  Estas classificações são \textbf{independentes}.

  \begin{itemize}

  \item
    Python é dinamicamente tipado e fortemente tipado; C é estaticamente
    tipado e fracamente tipado.
  \end{itemize}
\end{itemize}
\end{frame}

\hypertarget{a-linguagem-de-expressuxf5es-tipadas}{%
\section{A Linguagem de Expressões
Tipadas}\label{a-linguagem-de-expressuxf5es-tipadas}}

\begin{frame}{Sintaxe}
\protect\hypertarget{sintaxe}{}
\[\begin{array}{lcll}
t & ::= & \mathbf{true} \mid \mathbf{false} & \text{(booleanos)} \\
  & \mid & \mathbf{if}\ t_1\ \mathbf{then}\ t_2\ \mathbf{else}\ t_3 & \text{(condicional)} \\
  & \mid & \mathbf{0} & \text{(zero)} \\
  & \mid & \mathbf{succ}\ t \mid \mathbf{pred}\ t & \text{(sucessor, predecessor)} \\
  & \mid & \mathbf{iszero}\ t & \text{(teste de zero)}
\end{array}\]
\end{frame}

\begin{frame}{Tipos}
\protect\hypertarget{tipos}{}
\textbf{Tipos:} \[T ::= \mathbf{Bool} \mid \mathbf{Nat}\]
\end{frame}

\begin{frame}{Valores}
\protect\hypertarget{valores}{}
\begin{itemize}

\item
  \textbf{Valores} (termos que não podem ser reduzidos):

  \begin{itemize}

  \item
    \(nv\) são os \textbf{valores numéricos}: \(\mathbf{0}\),
    \(\mathbf{succ}\ \mathbf{0}\),
    \(\mathbf{succ}\ (\mathbf{succ}\ \mathbf{0})\), \ldots{}
  \item
    Valores booleanos: \(\mathbf{true}\) e \(\mathbf{false}\)
  \end{itemize}
\end{itemize}

\[\begin{array}{lcl}
v  & ::= & \mathbf{true} \mid \mathbf{false} \mid nv \\
nv & ::= & \mathbf{0} \mid \mathbf{succ}\ nv
\end{array}\]
\end{frame}



\begin{frame}{Relação de Tipagem}
\protect\hypertarget{relauxe7uxe3o-de-tipagem}{}
\begin{itemize}

\item
  Notação \(\vdash t : T\)

  \begin{itemize}

  \item
    o termo \(t\) tem tipo \(T\).
  \end{itemize}
\item
  A relação é definida por regras de inferência.
\end{itemize}
\end{frame}

\begin{frame}{Booleanos}
\protect\hypertarget{booleanos}{}
\[
\begin{array}{c}
\frac{}{\vdash \mathbf{true} : \mathbf{Bool}} \quad\text{(T-True)}\\
\frac{}{\vdash \mathbf{false} : \mathbf{Bool}} \quad\text{(T-False)}
\end{array}
\]
\end{frame}

\begin{frame}{Condicionais}
\protect\hypertarget{condicionais}{}
\[\frac{
  \vdash t_1 : \mathbf{Bool} \quad
  \vdash t_2 : T \quad
  \vdash t_3 : T
}{
  \vdash \mathbf{if}\ t_1\ \mathbf{then}\ t_2\ \mathbf{else}\ t_3 : T
} \quad\text{(T-If)}\]

A regra (T-If) exige que os dois ramos tenham o \textbf{mesmo tipo}
\(T\).
\end{frame}

\begin{frame}{Naturais}
\protect\hypertarget{naturais}{}
\[
\begin{array}{c}
   \frac{}{\vdash \mathbf{0} : \mathbf{Nat}} \quad\text{(T-Zero)}\\
   \frac{\vdash t_1 : \mathbf{Nat}}{\vdash \mathbf{succ}\ t_1 : \mathbf{Nat}} \quad\text{(T-Succ)}\\
\end{array}
\]
\end{frame}

\begin{frame}{Naturais}
\protect\hypertarget{naturais-1}{}
\[
\begin{array}{c}
\frac{\vdash t_1 : \mathbf{Nat}}{\vdash \mathbf{pred}\ t_1 : \mathbf{Nat}} \quad\text{(T-Pred)}\\
\frac{\vdash t_1 : \mathbf{Nat}}{\vdash \mathbf{iszero}\ t_1 : \mathbf{Bool}} \quad\text{(T-IsZero)}
\end{array}\]
\end{frame}

\begin{frame}{Exemplo}
\protect\hypertarget{exemplo}{}
\(\vdash \mathbf{iszero}\ (\mathbf{succ}\ \mathbf{0}) : \mathbf{Bool}\)

\[\dfrac{
  \dfrac{
    \dfrac{}{\vdash \mathbf{0} : \mathbf{Nat}}
  }{\vdash \mathbf{succ}\ \mathbf{0} : \mathbf{Nat}}
}{\vdash \mathbf{iszero}\ (\mathbf{succ}\ \mathbf{0}) : \mathbf{Bool}}\]
\end{frame}

\begin{frame}{Exemplo}
\protect\hypertarget{exemplo-1}{}
\begin{itemize}

\item
  \(\mathbf{succ}\ \mathbf{true}\) \textbf{não é tipável}:

  \begin{itemize}

  \item
    A regra (T-Succ) exige que o argumento tenha tipo \(\mathbf{Nat}\)
  \item
    Mas \(\mathbf{true}\) tem tipo \(\mathbf{Bool}\)
  \item
    Não existe derivação válida --- o tipo checker \textbf{rejeita} o
    programa
  \end{itemize}
\end{itemize}
\end{frame}

\begin{frame}{Exemplo}
\protect\hypertarget{exemplo-2}{}
\begin{itemize}

\item
  Outros exemplos de termos não tipáveis:

  \begin{itemize}

  \item
    \(\mathbf{pred}\ \mathbf{false}\): (T-Pred) exige \(\mathbf{Nat}\),
    argumento é \(\mathbf{Bool}\)
  \item
    \(\mathbf{if}\ \mathbf{0}\ \mathbf{then}\ t_2\ \mathbf{else}\ t_3\):
    (T-If) exige \(\mathbf{Bool}\) na condição
  \end{itemize}
\end{itemize}
\end{frame}

\begin{frame}{Unicidade de Tipos}
\protect\hypertarget{unicidade-de-tipos}{}
\begin{itemize}
\item
  \textbf{Lema (Unicidade):} Se \(\vdash t : T\) e \(\vdash t : T'\),
  então \(T = T'\).
\item
  A tipagem é \textbf{determinística}: o verificador sempre produz um
  resultado único.
\end{itemize}
\end{frame}

\hypertarget{semuxe2ntica-small-step}{%
\section{Semântica Small-Step}\label{semuxe2ntica-small-step}}

\begin{frame}{Formas Presas}
\protect\hypertarget{formas-presas}{}
\begin{itemize}

\item
  Os termos irredutíveis são \textbf{valores} \(v\) ou \textbf{formas
  presas} (\emph{stuck terms}):

  \begin{itemize}

  \item
    Termos que não são valores e para os quais nenhuma regra se aplica
  \item
    Representam \textbf{erros em tempo de execução}
  \end{itemize}
\end{itemize}
\end{frame}

\begin{frame}{Formas Presas}
\protect\hypertarget{formas-presas-1}{}
\begin{itemize}
\item
  Exemplos de formas presas:

  \begin{itemize}

  \item
    \(\mathbf{succ}\ \mathbf{true}\): \(\mathbf{true}\) não é \(nv\), e
    nenhuma regra reduz isso
  \item
    \(\mathbf{pred}\ \mathbf{false}\): idem
  \item
    \(\mathbf{if}\ \mathbf{0}\ \mathbf{then}\ t_2\ \mathbf{else}\ t_3\):
    \(\mathbf{0}\) é valor numérico, não booleano
  \end{itemize}
\end{itemize}
\end{frame}

\begin{frame}{Formas Presas}
\protect\hypertarget{formas-presas-2}{}
\begin{itemize}

\item
  \textbf{O papel do sistema de tipos}: garantir que termos bem tipados
  \textbf{nunca} se tornem formas presas.
\end{itemize}
\end{frame}

\begin{frame}{Regras de Redução}
\protect\hypertarget{regras-de-reduuxe7uxe3o}{}
\[\frac{t_1 \to t_1'}{\mathbf{if}\ t_1\ \mathbf{then}\ t_2\ \mathbf{else}\ t_3 \to \mathbf{if}\ t_1'\ \mathbf{then}\ t_2\ \mathbf{else}\ t_3}
\quad\text{(E-If)}\]
\end{frame}

\begin{frame}{Regras de Redução}
\protect\hypertarget{regras-de-reduuxe7uxe3o-1}{}
\[\frac{}{\mathbf{if}\ \mathbf{true}\ \mathbf{then}\ t_2\ \mathbf{else}\ t_3 \to t_2}
\quad\text{(E-IfTrue)}\]
\end{frame}

\begin{frame}{Regras de Redução}
\protect\hypertarget{regras-de-reduuxe7uxe3o-2}{}
\[\frac{}{\mathbf{if}\ \mathbf{false}\ \mathbf{then}\ t_2\ \mathbf{else}\ t_3 \to t_3}
\quad\text{(E-IfFalse)}\]
\end{frame}

\begin{frame}{Regras de Redução}
\protect\hypertarget{regras-de-reduuxe7uxe3o-3}{}
\[\frac{t_1 \to t_1'}{\mathbf{succ}\ t_1 \to \mathbf{succ}\ t_1'}
\quad\text{(E-Succ)}\]
\end{frame}

\begin{frame}{Regras de Redução}
\protect\hypertarget{regras-de-reduuxe7uxe3o-4}{}
\[
\frac{t_1 \to t_1'}{\mathbf{pred}\ t_1 \to \mathbf{pred}\ t_1'}
\quad\text{(E-Pred)}\]
\end{frame}

\begin{frame}{Regras de Redução}
\protect\hypertarget{regras-de-reduuxe7uxe3o-5}{}
\[\frac{}{\mathbf{pred}\ \mathbf{0} \to \mathbf{0}}
\quad\text{(E-PredZero)}
\]
\end{frame}

\begin{frame}{Regras de Redução}
\protect\hypertarget{regras-de-reduuxe7uxe3o-6}{}
\[
\frac{}{\mathbf{pred}\ (\mathbf{succ}\ nv_1) \to nv_1}
\quad\text{(E-PredSucc)}\]
\end{frame}

\begin{frame}{Regras de Redução}
\protect\hypertarget{regras-de-reduuxe7uxe3o-7}{}
\[\frac{t_1 \to t_1'}{\mathbf{iszero}\ t_1 \to \mathbf{iszero}\ t_1'}
\quad\text{(E-IsZero)}
\]
\end{frame}

\begin{frame}{Regras de Redução}
\protect\hypertarget{regras-de-reduuxe7uxe3o-8}{}
\[
\frac{}{\mathbf{iszero}\ \mathbf{0} \to \mathbf{true}}
\quad\text{(E-IsZeroZero)}\]
\end{frame}

\begin{frame}{Regras de Redução}
\protect\hypertarget{regras-de-reduuxe7uxe3o-9}{}
\[\frac{}{\mathbf{iszero}\ (\mathbf{succ}\ nv_1) \to \mathbf{false}}
\quad\text{(E-IsZeroSucc)}\]
\end{frame}

\begin{frame}{Exemplo}
\protect\hypertarget{exemplo-3}{}
\[\mathbf{if}\ (\mathbf{iszero}\ \mathbf{0})\ \mathbf{then}\ (\mathbf{succ}\ \mathbf{0})\ \mathbf{else}\ \mathbf{0}\]
\end{frame}

\begin{frame}{Exemplo}
\protect\hypertarget{exemplo-4}{}
\[\xrightarrow{\text{E-If, E-IsZeroZero}}
\mathbf{if}\ \mathbf{true}\ \mathbf{then}\ (\mathbf{succ}\ \mathbf{0})\ \mathbf{else}\ \mathbf{0}\]
\end{frame}

\begin{frame}{Exemplo}
\protect\hypertarget{exemplo-5}{}
\[\xrightarrow{\text{E-IfTrue}}
\mathbf{succ}\ \mathbf{0}\]
\end{frame}

\begin{frame}{Exemplo}
\protect\hypertarget{exemplo-6}{}
Outro exemplo:
\(\mathbf{pred}\ (\mathbf{succ}\ (\mathbf{succ}\ \mathbf{0})) \xrightarrow{\text{E-PredSucc}} \mathbf{succ}\ \mathbf{0}\)
\end{frame}

\hypertarget{corretude-do-sistema-de-tipos}{%
\section{Corretude do Sistema de
Tipos}\label{corretude-do-sistema-de-tipos}}

\begin{frame}{Type Soundness}
\protect\hypertarget{type-soundness}{}
\begin{quote}
\textbf{Teorema (Type Soundness):} Se \(\vdash t : T\) e \(t \to^* t'\),
então \(t'\) é um valor ou existe \(t''\) tal que \(t' \to t''\).
\end{quote}

Em palavras: \textbf{programas bem tipados não ficam presos}.
\end{frame}

\begin{frame}{Type Soundness}
\protect\hypertarget{type-soundness-1}{}
A prova apoia-se em dois lemas fundamentais e independentes:

\begin{table}
\begin{tabular}[]{@{}ll@{}}
\toprule
Lema & Enunciado \\
\midrule

\textbf{Preservação} & O tipo se mantém após cada passo de redução \\
\textbf{Progresso} & Um termo bem tipado é valor ou pode reduzir \\
\bottomrule
\end{tabular}
\end{table}
\end{frame}

\begin{frame}{Lema de Preservação}
\protect\hypertarget{lema-de-preservauxe7uxe3o}{}
\begin{quote}
\textbf{Lema (Preservation):} Se \(\vdash t : T\) e \(t \to t'\), então
\(\vdash t' : T\).
\end{quote}

\textbf{Significado:} a tipagem é \textbf{invariante sob redução} --- o
tipo não se altera durante a execução.
\end{frame}

\begin{frame}{Lema de Preservação}
\protect\hypertarget{lema-de-preservauxe7uxe3o-1}{}
\textbf{Por que é importante?} Sem preservação, um termo poderia começar
bem tipado e, após reduções, tornar-se de tipo diferente --- invalidando
qualquer garantia do verificador de tipos.
\end{frame}

\begin{frame}{Lema de Preservação}
\protect\hypertarget{lema-de-preservauxe7uxe3o-2}{}
\textbf{Prova:} Por indução na derivação de \(t \to t'\), analisando
cada regra de redução.
\end{frame}

\begin{frame}{Alguns Casos}
\protect\hypertarget{alguns-casos}{}
\emph{Caso} (E-IfTrue):
\(t = \mathbf{if}\ \mathbf{true}\ \mathbf{then}\ t_2\ \mathbf{else}\ t_3\),
\(t' = t_2\).

De (T-If): \(\vdash t_2 : T\) e \(\vdash t_3 : T\). Logo
\(\vdash t_2 : T\). ✓
\end{frame}

\begin{frame}{Alguns Casos}
\protect\hypertarget{alguns-casos-1}{}
\emph{Caso} (E-If): \(t_1 \to t_1'\),
\(t' = \mathbf{if}\ t_1'\ \mathbf{then}\ t_2\ \mathbf{else}\ t_3\).

De (T-If): \(\vdash t_1 : \mathbf{Bool}\). Por hip. de indução:
\(\vdash t_1' : \mathbf{Bool}\). Aplicando (T-If): \(\vdash t' : T\). ✓
\end{frame}

\begin{frame}{Alguns Casos}
\protect\hypertarget{alguns-casos-2}{}
\emph{Caso} (E-Succ): \(t_1 \to t_1'\), \(t' = \mathbf{succ}\ t_1'\).

De (T-Succ): \(\vdash t_1 : \mathbf{Nat}\). Por hip. de ind.:
\(\vdash t_1' : \mathbf{Nat}\). Logo
\(\vdash \mathbf{succ}\ t_1' : \mathbf{Nat}\). ✓
\end{frame}

\begin{frame}{Alguns Casos}
\protect\hypertarget{alguns-casos-3}{}
\emph{Caso} (E-IsZeroZero): \(t' = \mathbf{true}\). Por (T-True):
\(\vdash \mathbf{true} : \mathbf{Bool}\); tipo de \(t\) via (T-IsZero) é
\(\mathbf{Bool}\). ✓
\end{frame}

\begin{frame}{Lema de Progresso}
\protect\hypertarget{lema-de-progresso}{}
\begin{quote}
\textbf{Lema (Progress):} Se \(\vdash t : T\), então \(t\) é um valor ou
existe \(t'\) tal que \(t \to t'\).
\end{quote}

\textbf{Significado:} termos bem tipados \textbf{nunca ficam presos} ---
sempre progridem.
\end{frame}

\begin{frame}{Lema de Progresso}
\protect\hypertarget{lema-de-progresso-1}{}
\textbf{Por que é importante?} Sem progresso, o tipo seria preservado
mas a execução poderia travar. O progresso garante que nunca há um
estado de erro irrecuperável.

\textbf{Prova:} Por indução estrutural na derivação de \(\vdash t : T\).
\end{frame}

\begin{frame}{Alguns Casos}
\protect\hypertarget{alguns-casos-4}{}
\emph{Casos} (T-True), (T-False), (T-Zero): \(t\) já é valor. ✓
\end{frame}

\begin{frame}{Alguns Casos}
\protect\hypertarget{alguns-casos-5}{}
\emph{Caso} (T-If): \(\vdash t_1 : \mathbf{Bool}\). Por hip. de ind.
sobre \(t_1\): ou \(t_1\) é valor, ou reduz.

\begin{itemize}

\item
  Se \(t_1 \to t_1'\): aplica-se (E-If). ✓
\item
  Se \(t_1\) é valor de tipo \(\mathbf{Bool}\): pelo lema de unicidade,
  \(t_1\) não pode ser \(nv\) (tipo \(\mathbf{Nat}\)). Logo
  \(t_1 \in \{\mathbf{true}, \mathbf{false}\}\).

  \begin{itemize}

  \item
    \(t_1 = \mathbf{true}\): aplica-se (E-IfTrue). ✓
  \item
    \(t_1 = \mathbf{false}\): aplica-se (E-IfFalse). ✓
  \end{itemize}
\end{itemize}
\end{frame}

\begin{frame}{Alguns Casos}
\protect\hypertarget{alguns-casos-6}{}
\emph{Caso} (T-Succ): \(\vdash t_1 : \mathbf{Nat}\). Por hip. de ind.:
ou \(t_1\) reduz, ou é valor.

\begin{itemize}

\item
  \(t_1 \to t_1'\): aplica-se (E-Succ). ✓
\item
  \(t_1\) é valor de tipo \(\mathbf{Nat}\): pela unicidade, \(t_1\) é
  \(nv\). Logo \(\mathbf{succ}\ nv\) é valor numérico. ✓
\end{itemize}
\end{frame}

\begin{frame}{Alguns Casos}
\protect\hypertarget{alguns-casos-7}{}
\emph{Caso} (T-Pred): \(\vdash t_1 : \mathbf{Nat}\). Por hip. de ind.:
ou \(t_1\) reduz, ou é \(nv\).

\begin{itemize}

\item
  \(t_1 \to t_1'\): aplica-se (E-Pred). ✓ \textbar{}
  \(nv = \mathbf{0}\): (E-PredZero). ✓ \textbar{}
  \(nv = \mathbf{succ}\ nv_1\): (E-PredSucc). ✓
\end{itemize}
\end{frame}

\begin{frame}{Corretude}
\protect\hypertarget{corretude}{}
\begin{quote}
\textbf{Prova do Teorema (Type Soundness):} Seja \(\vdash t : T\) e
\(t \to^* t'\).

Por \textbf{preservação} (aplicada repetidamente a cada passo da
sequência \(t \to^* t'\)): \(\vdash t' : T\).

Por \textbf{progresso}: \(t'\) é valor ou existe \(t''\) com
\(t' \to t''\). ✓
\end{quote}
\end{frame}

\hypertarget{sistema-de-tipos-para-tline}{%
\section{Sistema de Tipos para
TLine}\label{sistema-de-tipos-para-tline}}

\begin{frame}{TLine: Contexto e Relação de Tipagem}
\protect\hypertarget{tline-contexto-e-relauxe7uxe3o-de-tipagem}{}
\begin{itemize}

\item
  TLine tem três tipos (\(\mathbf{Int}\), \(\mathbf{Bool}\),
  \(\mathbf{String}\)) e um \textbf{contexto de tipagem}:
\end{itemize}

\[\Gamma : \mathit{Var} \rightharpoonup \mathit{Ty}\]
\end{frame}

\begin{frame}{Regras de Tipos}
\protect\hypertarget{regras-de-tipos}{}
\begin{itemize}

\item
  Regras para expressões: \(\Gamma \vdash e : T\)
\end{itemize}

\[\frac{}{\Gamma \vdash n : \mathbf{Int}} \quad\text{(T-Int)}\]
\end{frame}

\begin{frame}{Regras de Tipos}
\protect\hypertarget{regras-de-tipos-1}{}
\[
\frac{x \in \mathrm{dom}(\Gamma)}{\Gamma \vdash x : \Gamma(x)} \quad\text{(T-Var)}\]
\end{frame}

\begin{frame}{Regras de Tipos}
\protect\hypertarget{regras-de-tipos-2}{}
\[\frac{\Gamma \vdash e_1 : \mathbf{Int} \quad \Gamma \vdash e_2 : \mathbf{Int}}{\Gamma \vdash e_1 \oplus e_2 : \mathbf{Int}} \quad\text{(T-ArithOp)}
\]
\end{frame}

\begin{frame}{Regras de Tipos}
\protect\hypertarget{regras-de-tipos-3}{}
\[
\frac{\Gamma \vdash e_1 : \mathbf{Int} \quad \Gamma \vdash e_2 : \mathbf{Int}}{\Gamma \vdash e_1 \otimes e_2 : \mathbf{Bool}} \quad\text{(T-CmpOp)}\]
\end{frame}

\begin{frame}{Regras de Tipos}
\protect\hypertarget{regras-de-tipos-4}{}
\begin{itemize}

\item
  Regras para comandos: \(\Gamma \vdash s \dashv \Gamma'\)
\end{itemize}

\[\frac{\Gamma \vdash e : T}{\Gamma \vdash \mathbf{var}\ x : T = e \dashv \Gamma[x \mapsto T]} \quad\text{(T-Decl)}\]
\end{frame}

\begin{frame}{Regras de Tipos}
\protect\hypertarget{regras-de-tipos-5}{}
\[\frac{x \in \mathrm{dom}(\Gamma) \quad \Gamma \vdash e : \Gamma(x)}{\Gamma \vdash x \mathbin{:=} e \dashv \Gamma} \quad\text{(T-Assign)}\]
\end{frame}

\begin{frame}{Regras de Tipos}
\protect\hypertarget{regras-de-tipos-6}{}
\[\frac{\Gamma \vdash e_p : \mathbf{String} \quad x \in \mathrm{dom}(\Gamma) \quad \Gamma(x) = \mathbf{Int}}{\Gamma \vdash \mathbf{read}\ e_p\ x \dashv \Gamma} \quad\text{(T-Read)}\]
\end{frame}

\begin{frame}{Regras de Tipos}
\protect\hypertarget{regras-de-tipos-7}{}
\begin{itemize}

\item
  \textbf{T-Decl} é a única regra que \textbf{estende} o contexto:
  \(\Gamma' = \Gamma[x \mapsto T]\)
\item
  As demais \textbf{preservam} o contexto: \(\Gamma' = \Gamma\)
\end{itemize}
\end{frame}

\hypertarget{sistema-de-tipos-para-twhile}{%
\section{Sistema de Tipos para
TWhile}\label{sistema-de-tipos-para-twhile}}

\begin{frame}[fragile]{Escopo}
\protect\hypertarget{escopo}{}
\begin{itemize}

\item
  TWhile adiciona \texttt{if} e
  \texttt{while} com \textbf{escopo léxico}: variáveis
  declaradas dentro de um bloco não escapam.
\end{itemize}
\end{frame}

\begin{frame}{Escopo}
\protect\hypertarget{escopo-1}{}
\textbf{Formalização}: o contexto de \textbf{saída} é sempre igual ao
contexto de \textbf{entrada} \(\Gamma\):

\[\frac{\Gamma \vdash e : \mathbf{Bool} \quad \Gamma \vdash B_1 \dashv \Gamma_1 \quad \Gamma \vdash B_2 \dashv \Gamma_2}{\Gamma \vdash \mathbf{if}\ e\ \mathbf{then}\ B_1\ \mathbf{else}\ B_2 \dashv \Gamma} \quad\text{(T-If)}\]
\end{frame}

\begin{frame}{Escopo}
\protect\hypertarget{escopo-2}{}
\[\frac{\Gamma \vdash e : \mathbf{Bool} \quad \Gamma \vdash B \dashv \Gamma'}{\Gamma \vdash \mathbf{while}\ e\ \mathbf{do}\ B \dashv \Gamma} \quad\text{(T-While)}\]

\begin{itemize}

\item
  Os contextos \(\Gamma_1\), \(\Gamma_2\) e \(\Gamma'\) são
  \textbf{descartados}: variáveis locais aos blocos não são visíveis
  fora deles.
\end{itemize}
\end{frame}

\hypertarget{sistema-de-tipos-para-timp}{%
\section{Sistema de Tipos para TImp}\label{sistema-de-tipos-para-timp}}

\begin{frame}{Tipos e Ambientes}
\protect\hypertarget{tipos-e-ambientes}{}
\begin{itemize}

\item
  TImp estende TWhile com funções e registros
\end{itemize}

\[\tau \;::=\; \mathbf{Int} \mid \mathbf{Bool} \mid \mathbf{String} \mid R \qquad \rho \;::=\; \mathbf{void} \mid \tau\]
\end{frame}

\begin{frame}{Tipos e Ambientes}
\protect\hypertarget{tipos-e-ambientes-1}{}
\begin{itemize}

\item
  Relação de tipagem: \(\Gamma;\Delta;\Phi \vdash e : \tau\) e
  \(\Gamma;\Delta;\Phi;\rho_\mathit{ret} \vdash s \dashv \Gamma'\)
\end{itemize}
\end{frame}

\begin{frame}{Regras para Registros}
\protect\hypertarget{regras-para-registros}{}
\textbf{Acesso a campo (T-Field):}
\[\frac{\Gamma;\Delta;\Phi \vdash e : R \quad \Delta(R)(f) = \tau}{\Gamma;\Delta;\Phi \vdash e.f : \tau}\]
\end{frame}

\begin{frame}{Regras para Registros}
\protect\hypertarget{regras-para-registros-1}{}
\textbf{Construção (T-New):}
\[\frac{R \in \mathrm{dom}(\Delta) \quad \Delta(R) = \overline{(f_i:\tau_i)} \quad \Gamma;\Delta;\Phi \vdash e_i : \tau_i}{\Gamma;\Delta;\Phi \vdash \mathbf{new}\ R\{\overline{f_i=e_i}\} : R}\]
\end{frame}

\begin{frame}{Regras para Registros}
\protect\hypertarget{regras-para-registros-2}{}
\textbf{Atribuição de campo (T-FieldAssign):}
\[\frac{\Gamma(x) = R \quad \Delta(R)(f) = \tau \quad \Gamma;\Delta;\Phi \vdash e : \tau}{\Gamma;\Delta;\Phi \vdash x.f := e \dashv \Gamma}\]
\end{frame}

\begin{frame}{Regras para Funções}
\protect\hypertarget{regras-para-funuxe7uxf5es}{}
\textbf{Chamada como expressão (T-CallExpr):}
\[\frac{\Phi(f) = ([\tau_1,\ldots,\tau_n],\;\tau) \quad \Gamma;\Delta;\Phi \vdash e_i : \tau_i}{\Gamma;\Delta;\Phi \vdash f(e_1,\ldots,e_n) : \tau}\]
\end{frame}

\begin{frame}{Regras para Funções}
\protect\hypertarget{regras-para-funuxe7uxf5es-1}{}
\textbf{Retorno com valor (T-RetVal):}
\[\frac{\rho_\mathit{ret} = \tau \quad \Gamma;\Delta;\Phi \vdash e : \tau}{\Gamma;\ldots;\rho_\mathit{ret} \vdash \mathbf{return}\ e \dashv \Gamma}\]
\end{frame}

\begin{frame}{Regras para Funções}
\protect\hypertarget{regras-para-funuxe7uxf5es-2}{}
\textbf{Declaração de função (T-FuncDecl):}
\[\frac{[x_1\!\mapsto\!\tau_1,\ldots,x_n\!\mapsto\!\tau_n];\Delta;\Phi;\rho \vdash B \dashv \_}{\Delta;\Phi \vdash \mathbf{fn}\ f(\overline{x_i:\tau_i}):\rho\ \{B\}\ \checkmark}\]
\end{frame}

\hypertarget{conclusuxe3o}{%
\section{Conclusão}\label{conclusuxe3o}}

\begin{frame}{Conclusão}
\protect\hypertarget{conclusuxe3o-1}{}
\begin{itemize}

\item
  \textbf{Tipos} classificam valores e restringem operações; tipagem
  pode ser estática ou dinâmica, forte ou fraca
\item
  O \textbf{sistema de tipos} \(\vdash t : T\) é definido por regras de
  inferência --- cada construtor tem exatamente uma regra (unicidade de
  tipos)
\end{itemize}
\end{frame}

\begin{frame}{Conclusão}
\protect\hypertarget{conclusuxe3o-2}{}
\begin{itemize}

\item
  \textbf{Formas presas}: termos sem valor e sem redução possível ---
  representam erros em tempo de execução
\item
  \textbf{Type Soundness} = Preservação + Progresso: programas bem
  tipados nunca ficam presos
\end{itemize}
\end{frame}

\begin{frame}{Conclusão}
\protect\hypertarget{conclusuxe3o-3}{}
\begin{itemize}

\item
  \textbf{Preservação}: prova que o tipo não muda durante a execução
  (indução na redução \(t \to t'\))
\item
  \textbf{Progresso}: prova que a execução nunca trava (indução na
  tipagem \(\vdash t : T\))
\item
  Os dois lemas são \textbf{independentes} e \textbf{complementares} ---
  cada um é necessário e insuficiente sozinho
\end{itemize}
\end{frame}
